% Curriculum Vitae -- Eonyong Han
% Build with latexmk -pdf cv.tex, or tectonic cv.tex.
\documentclass[11pt,a4paper]{article}
\usepackage[T1]{fontenc}
\usepackage[utf8]{inputenc}
\usepackage{lmodern}
\usepackage[margin=2cm]{geometry}
\usepackage{enumitem,titlesec,tabularx,xcolor,fancyhdr}
\usepackage{hyperref}
\definecolor{accent}{HTML}{1A3A5A}
\hypersetup{colorlinks=true,linkcolor=accent,urlcolor=accent,
  pdftitle={Curriculum Vitae -- Eonyong Han},pdfauthor={Eonyong Han}}
\titleformat{\section}{\large\bfseries\color{accent}}{}{0pt}{}[\vspace{-4pt}\hrule]
\titlespacing{\section}{0pt}{11pt}{7pt}
\setlength{\parindent}{0pt}
\setlength{\parskip}{3pt}
\setlength{\emergencystretch}{2em}
\setlist[itemize]{leftmargin=1.3em,itemsep=2pt,topsep=3pt,parsep=0pt}
\setlist[enumerate]{leftmargin=1.6em,itemsep=7pt,topsep=4pt,parsep=0pt}
\newcommand{\me}[1]{\textbf{#1}}
\newcommand{\entry}[2]{\begin{tabularx}{\linewidth}{@{}>{\raggedright\arraybackslash}X r@{}}\textbf{#1}&\small\color{accent}#2\end{tabularx}\par}
\newcommand{\subentry}[1]{{\itshape #1}\par}
\newcommand{\project}[4]{\vspace{3pt}\entry{#1}{#2}{\small #3\enspace\textbullet\enspace PI: #4}\par}
\pagestyle{fancy}
\fancyhf{}
\renewcommand{\headrulewidth}{0pt}
\fancyfoot[L]{\footnotesize\textcolor{accent}{Eonyong Han\enspace |\enspace Curriculum Vitae}}
\fancyfoot[R]{\footnotesize September 2026\enspace |\enspace\thepage}

\begin{document}
\begin{center}
  {\Huge\bfseries\color{accent}Eonyong Han}\\[5pt]
  {\large Ph.D.\ Candidate, Bioinformatics \& Computational Biology}\\[3pt]
  School of Computer Science and Engineering, Kyungpook National University\\
  Daegu 41566, Republic of Korea\\[5pt]
  \href{mailto:djssoddl@knu.ac.kr}{djssoddl@knu.ac.kr}\quad\textbullet\quad
  \href{https://scholar.google.com/citations?user=NMRL9agAAAAJ}{Google Scholar}\quad\textbullet\quad
  \href{https://github.com/EONYONG-HAN}{GitHub}\quad\textbullet\quad
  \href{https://orcid.org/0000-0002-2426-4962}{ORCID: 0000-0002-2426-4962}
\end{center}

\section{Research Profile}
I develop computational methods for transcriptome reconstruction, multi-omics integration,
and immune-response interpretation. My dissertation connects computational performance
with biological interpretation by examining how data selection and partitioning shape
the evidence available for inference. My background combines computer science with
experimental protein research.

\section{Education}
\entry{Ph.D., Computer Science and Engineering}{Sep.\ 2022 -- Feb.\ 2027 (expected)}
Kyungpook National University\enspace\textbullet\enspace Advisor: Prof.\ Inuk Jung\par
Dissertation: \emph{Data Selection and Partitioning for Transcriptome Reconstruction,
Multi-omics Integration, and Immune-response Interpretation}
\vspace{4pt}

\entry{M.S., Computer Science and Engineering}{2020 -- 2022}
Kyungpook National University\par
Thesis: \emph{A Study on Network Comparison with Invariant Topology Size Using Colon Cancer Data}
\vspace{4pt}

\entry{B.S., Biotechnology}{2010 -- 2019}
School of Life Science and Biotechnology, Kyungpook National University

\section{Research Experience}
\entry{Graduate Researcher}{2020 -- Present}
\subentry{Computational Biology \& Bioinformatics Lab, Kyungpook National University}
Advisor: Prof.\ Inuk Jung
\begin{itemize}
  \item Developed Deep-SemP, a semantic read-partitioning framework for parallel
    transcriptome assembly, using protein-derived representations and knowledge
    distillation to a nucleotide model.
  \item Conducted de novo transcriptome reconstruction and gene-expression analysis
    in \emph{Cnidium officinale} under heat stress.
  \item Benchmarked ten clustering methods for multi-omics study design; developed
    pathway and network analyses for TCGA cancer data.
  \item Built random-forest and SHAP analyses of COVID-19 immune response; contributed
    to density-based analysis of SARS-CoV-2 mutation hotspots.
\end{itemize}
\vspace{3pt}
\entry{Research Intern, then Graduate Researcher}{Dec.\ 2018 -- Aug.\ 2020}
\subentry{Structural Biology, School of Life Science and Biotechnology, KNU}
Supervisor: Prof.\ Eric Di Luccio
\begin{itemize}
  \item Produced recombinant proteins in \emph{E. coli}; performed SDS-PAGE, western
    blotting, chromatographic purification, and vapour-diffusion crystallization.
  \item Contributed to an NRF-funded study of Set7 methyltransferase in
    \emph{Schizosaccharomyces pombe} gametogenesis and its association with human infertility.
\end{itemize}

\newpage
\section{Publications}
{\small $^{\dagger}$Equal contribution. Author's name in \textbf{bold}.}

\subentry{Under review}
\begin{enumerate}
  \item \me{Han, E.}$^{\dagger}$, Park, J.$^{\dagger}$, Kim, Y., \& Jung, I. (2026).
    A deep semantic read-partitioning framework for parallel reference-free de novo
    transcriptome assembly in \emph{Caenorhabditis}.
    Under review at \emph{Bioinformatics}.\\
    Code: \href{https://github.com/cobi-git/Deep-SemP}{github.com/cobi-git/Deep-SemP}
\end{enumerate}

\subentry{Peer-reviewed journal articles}
\begin{enumerate}
  \item Shin, S.$^{\dagger}$, \me{Han, E.}$^{\dagger}$, Seong, H., Kim, Y.I., Jung, I., \& Jung, W. (2025).
    De novo transcriptome assembly and gene expression analysis of \emph{Cnidium officinale}
    under high-temperature conditions. \emph{BMC Genomics}, 26, 907.\\
    \href{https://doi.org/10.1186/s12864-025-12051-5}{doi:10.1186/s12864-025-12051-5}
  \item \me{Han, E.}$^{\dagger}$, Kwon, H.$^{\dagger}$, \& Jung, I. (2025).
    A review on multi-omics integration for aiding study design of large scale TCGA
    cancer datasets. \emph{BMC Genomics}, 26, 769.\\
    \href{https://doi.org/10.1186/s12864-025-11925-y}{doi:10.1186/s12864-025-11925-y}
  \item Youn, S.$^{\dagger}$, Jeong, D.$^{\dagger}$, Kwon, H.$^{\dagger}$, \me{Han, E.}, Kim, S., \& Jung, I. (2025).
    Identification of severity related mutation hotspots in SARS-CoV-2 using a
    density-based clustering approach. \emph{BioData Mining}, 18, 61.\\
    \href{https://doi.org/10.1186/s13040-025-00476-3}{doi:10.1186/s13040-025-00476-3}
  \item \me{Han, E.}, Youn, S., Kwon, K.T., Kim, S.C., Jo, H.-Y., \& Jung, I. (2024).
    Disease progression associated cytokines in COVID-19 patients with deteriorating
    and recovering health conditions. \emph{Scientific Reports}, 14, 24712.\\
    \href{https://doi.org/10.1038/s41598-024-75924-x}{doi:10.1038/s41598-024-75924-x}
  \item Jeon, J., \me{Han, E.Y.}, \& Jung, I. (2023).
    MOPA: An integrative multi-omics pathway analysis method for measuring omics
    activity. \emph{PLOS ONE}, 18(3), e0278272.\\
    \href{https://doi.org/10.1371/journal.pone.0278272}{doi:10.1371/journal.pone.0278272}
\end{enumerate}

\section{Technical Skills}
\begin{itemize}
  \item \textbf{Programming and tools:} Python, R, Bash, Git, LaTeX, Linux/HPC.
  \item \textbf{Genomics:} RNA-seq and differential expression, de novo transcriptome
    assembly, single-cell RNA-seq, variant/mutation analysis, miRNA--target analysis.
  \item \textbf{Pathways and networks:} KEGG/GO enrichment, WGCNA, pathway-unit
    correlation networks, graph edit distance, network propagation.
  \item \textbf{Machine learning:} Random forests, SHAP, density-based and consensus
    clustering, feature selection, benchmarking design.
  \item \textbf{Deep learning:} ESM-2, DNABERT-2, knowledge distillation, transformer
    fine-tuning, multi-GPU inference with PyTorch.
  \item \textbf{Assembly and computing:} Trinity, BUSCO, minimap2, StringTie, CD-HIT-EST,
    CAP3, Jellyfish; complexity-aware CPU scheduling and multi-server parallelization.
  \item \textbf{Data resources:} TCGA, GISAID, NCBI SRA, Ensembl, UniProt, KEGG;
    clinical EHR and cytokine-profile cohorts.
\end{itemize}

\newpage
\section{Funded Research Projects}
Selected participation as a student researcher or research assistant.
Principal investigators (PIs) are identified below.

\project{Petascale in-storage computing for accelerating genomic DNA sequence analysis}
  {Sep.\ 2024 -- Jun.\ 2026}{NRF, Basic Research Laboratory}{Yongtae Kim}
\project{Quantum advantage for quantum-computing-based genomic DNA alignment}
  {Jun.\ 2025 -- Dec.\ 2025}{NRF}{Yongtae Kim}
\project{Time-series multi-omics integrative analysis and AI applications}
  {Apr.\ 2024 -- Dec.\ 2025}{KDCA}{Inuk Jung}
\project{Multi-omics integrative analysis and prognosis prediction modelling for COVID-19 patients}
  {Mar.\ 2022 -- Dec.\ 2023}{KDCA}{Inuk Jung}
\project{Development of an AI-based infection severity classification system}
  {Aug.\ 2022 -- Dec.\ 2024}{KHIDI}{Inuk Jung}
\project{Large-scale multi-omics data integration methods and open platform}
  {Mar.\ 2021 -- Dec.\ 2021}{NRF}{Inuk Jung}
\project{Machine learning methods for single-cell time-series transcriptome data}
  {Mar.\ 2021 -- May 2022}{NRF}{Inuk Jung}
\project{AI diagnostic technology for embedded imaging medical devices}
  {May 2022 -- Dec.\ 2024}{KEIT}{Jaeil Kim}
\project{High-performance spatiotemporal IoT data management and intelligence engine}
  {Mar.\ 2024 -- May 2025}{NRF}{Youngkyoon Seo}
\project{Usage-pattern classification for personalized air-conditioner operation}
  {May 2022 -- Apr.\ 2023}{LG Electronics, industry collaboration}{Inuk Jung}
\project{Role of Set7 in \emph{Schizosaccharomyces pombe} gametogenesis and its association with human infertility}
  {Sep.\ 2019 -- Aug.\ 2020}{NRF}{Eric Di Luccio}

\section{Fellowship}
\entry{BK21 FOUR}{Sep.\ 2020 -- Aug.\ 2026}
Intelligence-Convergence Software Education \& Research Group\par
National Research Foundation of Korea\enspace\textbullet\enspace Participating graduate student

\section{Reference}
\entry{Prof.\ Inuk Jung}{Ph.D.\ Advisor}
School of Computer Science and Engineering, Kyungpook National University\par
\href{mailto:inukjung@knu.ac.kr}{inukjung@knu.ac.kr}
\end{document}
